\documentclass[11pt,a4paper]{article}

% ---------- Packages ----------
\usepackage[utf8]{inputenc}
\usepackage[T1]{fontenc}
\usepackage{lmodern}
\usepackage[french]{babel}
\usepackage{amsmath,amssymb}
\usepackage[margin=2.3cm,headheight=15pt]{geometry}
\usepackage{fancyhdr}
\usepackage{enumitem}
\usepackage{xcolor}
\usepackage[most]{tcolorbox}
\usepackage{hyperref}
\usepackage{tikz}

% ---------- En-tête / pied de page ----------
\pagestyle{fancy}
\fancyhf{}
\fancyhead[L]{\small Gauthier Baudelet - 0472522463}
\fancyhead[R]{\small Trigonométrie : Les bases}
\fancyfoot[L]{\tiny Usage strictement personnel -- Diffusion interdite}
\fancyfoot[C]{\small\thepage}
\renewcommand{\headrulewidth}{0.4pt}

% ---------- Mise en forme ----------
\setlength{\parindent}{0pt}
\setlength{\parskip}{0.5em}
\setlist[itemize]{itemsep=2pt,topsep=2pt}
\setlist[enumerate]{itemsep=2pt,topsep=2pt}

\newtcolorbox{definition}{breakable,colback=blue!4,colframe=blue!60!black,
  fonttitle=\bfseries,title={Définition}}
\newtcolorbox{propriete}[1][Propriétés]{breakable,colback=green!5!white,colframe=green!45!black,
  fonttitle=\bfseries,title={#1}}
\newtcolorbox{attention}{breakable,colback=red!4,colframe=red!60!black,
  fonttitle=\bfseries,title={Attention}}
\newtcolorbox{methode}{breakable,colback=orange!6,colframe=orange!70!black,
  fonttitle=\bfseries,title={Méthode}}

\newcommand{\dec}{{,}} % virgule décimale en mode math

\title{\textbf{Trigonométrie dans le triangle rectangle}\\[2pt]\large Synthèse de cours et exercices (3e secondaire)}
\author{}
\date{}

\begin{document}
\maketitle
\thispagestyle{fancy}
\vspace{-2cm}
\textbf{Propriété intellectuelle :} Ce document est rédigé et mis à disposition exclusivement pour les élèves inscrits aux cours de Gauthier Baudelet. Toute reproduction, mise en ligne (Google Drive partagé, réseaux sociaux, plateformes publiques) ou transmission à des tiers, même à titre gratuit, est strictement interdite.
% =====================================================
\section{Le triangle rectangle et la terminologie}
% =====================================================

La trigonométrie (du grec *trigonon*, triangle, et *metron*, mesure) s'intéresse aux relations entre les longueurs des côtés et les angles d'un triangle. En 3e secondaire, nous débutons dans le \textbf{triangle rectangle} (qui possède un angle de $90^\circ$).

\subsection{Nommer les côtés par rapport à un angle}

Dans un triangle rectangle, les côtés portent des noms précis selon l'angle aigu $\alpha$ que l'on étudie :

\begin{center}
\begin{tikzpicture}[scale=1.2]
  % Sommets du triangle
  \coordinate (A) at (0,0);
  \coordinate (B) at (4,0);
  \coordinate (C) at (0,3);
  
  % Tracé du triangle
  \draw[thick, blue!60!black] (A) -- (B) -- (C) -- cycle;
  
  % Angle droit en A
  \draw[thick] (0.3,0) -- (0.3,0.3) -- (0,0.3);
  
  % Étiquettes des sommets
  \node[left] at (A) {$A$};
  \node[right] at (B) {$B$};
  \node[above] at (C) {$C$};
  
  % Étiquettes des côtés
  \node[below] at (2,0) {\textbf{Côté adjacent} à $\alpha$};
  \node[left] at (0,1.5) {\textbf{Côté opposé} à $\alpha$};
  \node[above right] at (2,1.5) {\textbf{Hypoténuse}};
  
  % Arc de l'angle alpha en B
  \draw [thick, orange!70!black] (3.3,0) arc (180:143.13:0.7);
  \node at (3.1,0.3) {$\alpha$};
\end{tikzpicture}
\end{center}

\begin{definition}
Soit un triangle rectangle.
\begin{itemize}
  \item L'\textbf{hypoténuse} est le côté le plus long. C'est toujours celui qui se trouve \textbf{en face de l'angle droit} ($90^\circ$). Il ne change jamais de place.
  \item Le \textbf{côté opposé} à l'angle $\alpha$ est le côté situé \textbf{en face} de cet angle $\alpha$.
  \item Le \textbf{côté adjacent} à l'angle $\alpha$ est le côté qui \textbf{touche} l'angle $\alpha$ (sans être l'hypoténuse).
\end{itemize}
\end{definition}

\begin{attention}
Le côté opposé et le côté adjacent dépendent de l'angle $\alpha$ que l'on regarde ! Si l'on change d'angle aigu dans le triangle, le rôle de ces deux côtés s'inverse.
\end{attention}

% =====================================================
\section{Les trois rapports trigonométriques}
% =====================================================

Dans un triangle rectangle, le rapport entre deux longueurs de côtés ne dépend que de la mesure de l'angle aigu, et non de la taille du triangle. On définit ainsi trois fonctions : le \textbf{sinus}, le \textbf{cosinus} et la \textbf{tangente}.

\begin{propriete}[Les formules fondamentales -- SOH CAH TOA]
Pour un angle aigu $\alpha$ dans un triangle rectangle, on a :
\[
\sin(\alpha) = \frac{\text{Côté opposé}}{\text{hypoténuse}} \qquad (\text{\textbf{SOH}})
\]
\[
\cos(\alpha) = \frac{\text{Côté adjacent}}{\text{hypoténuse}} \qquad (\text{\textbf{CAH}})
\]
\[
\tan(\alpha) = \frac{\text{Côté opposé}}{\text{Côté adjacent}} \qquad (\text{\textbf{TOA}})
\]
\end{propriete}

\textbf{Astuce mnémotechnique :} Retenez le mot \textbf{SOH-CAH-TOA} :
\begin{itemize}
  \item \textbf{S}in = \textbf{O}pposé / \textbf{H}ypothénuse
  \item \textbf{C}os = \textbf{A}djacent / \textbf{H}ypothénuse
  \item \textbf{T}an = \textbf{O}pposé / \textbf{A}djacent
\end{itemize}

% =====================================================
\section{Utilisations pratiques}
% =====================================================

\subsection{Calculer la longueur d'un côté}

\begin{methode}
Pour trouver un côté inconnu :
\begin{enumerate}
  \item Repérer l'angle connu et identifier les côtés impliqués (opposé, adjacent ou hypoténuse).
  \item Choisir le rapport trigonométrique adéquat ($\sin$, $\cos$ ou $\tan$).
  \item Poser l'équation et isoler l'inconnue à l'aide d'un produit en croix.
\end{enumerate}
\end{methode}

\textbf{Exemple :} Soit un triangle rectangle où l'angle $\alpha = 30^\circ$, l'hypoténuse vaut $10\text{ cm}$, et l'on cherche le côté opposé $x$.
\[
\sin(30^\circ) = \frac{x}{10} \quad\Longrightarrow\quad x = 10 \times \sin(30^\circ) = 10 \times 0\dec 5 = 5\text{ cm}.
\]

\subsection{Calculer la mesure d'un angle}

\begin{methode}
Lorsque l'on connaît deux côtés et que l'on cherche l'angle $\alpha$ :
\begin{enumerate}
  \item Écrire le rapport trigonométrique reliant les deux côtés connus.
  \item Utiliser la fonction inverse sur la calculatrice ($\sin^{-1}$, $\cos^{-1}$ ou $\tan^{-1}$, souvent activée via la touche \texttt{2nd} ou \texttt{Shift}).
\end{enumerate}
\end{methode}

\textbf{Exemple :} Si $\tan(\alpha) = \dfrac{3}{4} = 0\dec 75$, alors $\alpha = \tan^{-1}(0\dec 75) \approx 36\dec 87^\circ$.

% =====================================================
\section{Récapitulatif}
% =====================================================

\begin{tcolorbox}[colback=gray!6,colframe=gray!60!black,title={\bfseries À retenir}]
\[
\sin(\alpha) = \frac{\text{Opp}}{\text{Hyp}}, \qquad \cos(\alpha) = \frac{\text{Adj}}{\text{Hyp}}, \qquad \tan(\alpha) = \frac{\text{Opp}}{\text{Adj}}
\]
\textbf{Règles d'or :}
\begin{itemize}
  \item Vérifiez toujours que votre calculatrice est bien réglée en \textbf{Degrés (DEG)} et non en Radians.
  \item L'hypoténuse est \textbf{toujours} plus grande que les côtés opposé et adjacent (donc le sinus et le cosinus d'un angle aigu sont toujours compris entre $0$ et $1$).
\end{itemize}
\end{tcolorbox}

\newpage
% =====================================================
\section{Exercices}
% =====================================================

\subsection{Série A -- Identification et application directe}

\begin{enumerate}[label=\textbf{\arabic*.}]

\item Dans un triangle $ABC$ rectangle en $A$, écris les formules de $\sin(B)$, $\cos(B)$ et $\tan(B)$ en fonction des côtés du triangle.

\item Calcule la valeur arrondie au centième (vérifie bien le mode de ta calculatrice) :
\begin{enumerate}[label=\alph*)]
  \item $\sin(45^\circ)$ \qquad
  \item $\cos(60^\circ)$ \qquad
  \item $\tan(30^\circ)$ \qquad
  \item $\sin(72^\circ)$
\end{enumerate}

\item Trouve la mesure de l'angle $\alpha$ au dixième de degré près :
\begin{enumerate}[label=\alph*)]
  \item $\sin(\alpha) = 0\dec 5$ \qquad
  \item $\cos(\alpha) = 0\dec 8$ \qquad
  \item $\tan(\alpha) = 2$ \qquad
  \item $\sin(\alpha) = \dfrac{3}{5}$
\end{enumerate}

\end{enumerate}

\subsection{Série B -- Calculs de longueurs et d'angles}

\begin{enumerate}[label=\textbf{\arabic*.},start=4]

\item Soit un triangle $RST$ rectangle en $S$. On sait que l'angle $\widehat{T} = 40^\circ$ et que l'hypoténuse $RT = 12\text{ cm}$. Calcule la longueur du côté $ST$ (adjacent à $\widehat{T}$).

\item Soit un triangle $MNP$ rectangle en $N$. Le côté opposé à l'angle $M$ mesure $5\text{ cm}$ et le côté adjacent mesure $7\text{ cm}$. Calcule la mesure de l'angle $\widehat{M}$.

\item Une échelle de $6\text{ mètres}$ de long est apposée contre un mur vertical. Le pied de l'échelle est situé à $2\text{ mètres}$ du mur. À quel angle (par rapport au sol) l'échelle est-elle inclinée ?

\end{enumerate}

\subsection{Défi bonus}

\begin{enumerate}[label=\alph*)]
  \item Démontre que pour tout angle aigu $\alpha$ d'un triangle rectangle, $\tan(\alpha) = \dfrac{\sin(\alpha)}{\cos(\alpha)}$.
  \item Sachant que $\sin^2(\alpha) + \cos^2(\alpha) = 1$ (relation fondamentale de trigonométrie), calcule $\cos(\alpha)$ si $\sin(\alpha) = 0\dec 6$.
\end{enumerate}

\newpage
% =====================================================
\section{Corrigés}
% =====================================================

\subsection*{Série A}

\textbf{1.} $\sin(B) = \dfrac{\text{opposé}}{\text{hypoténuse}}$, $\cos(B) = \dfrac{\text{adjacent}}{\text{hypoténuse}}$, $\tan(B) = \dfrac{\text{opposé}}{\text{adjacent}}$.

\textbf{2.} a) $0\dec 71$ \quad b) $0\dec 5$ \quad c) $0\dec 58$ \quad d) $0\dec 95$

\textbf{3.} a) $\alpha = 30^\circ$ \quad b) $\alpha \approx 36\dec 9^\circ$ \quad c) $\alpha \approx 63\dec 4^\circ$ \quad d) $\alpha \approx 36\dec 9^\circ$

\subsection*{Série B}

\textbf{4.} Dans le triangle $RST$, le côté $ST$ est adjacent à l'angle $\widehat{T}$ et on connaît l'hypoténuse $RT$. On utilise donc le cosinus :
\[
\cos(40^\circ) = \frac{ST}{12} \quad\Longrightarrow\quad ST = 12 \times \cos(40^\circ) \approx 12 \times 0\dec 766 = 9\dec 19\text{ cm}.
\]

\textbf{5.} On utilise la tangente :
\[
\tan(M) = \frac{\text{opposé}}{\text{adjacent}} = \frac{5}{7} \approx 0\dec 7143 \quad\Longrightarrow\quad M = \tan^{-1}(0\dec 7143) \approx 35\big[{^\circ}\big]5 \text{ ou } 35\dec 5^\circ.
\]

\textbf{6.} L'échelle forme l'hypoténuse ($6\text{ m}$) et la distance au sol forme le côté adjacent à l'angle cherché $\alpha$ entre le sol et l'échelle :
\[
\cos(\alpha) = \frac{2}{6} = \frac{1}{3} \quad\Longrightarrow\quad \alpha = \cos^{-1}\left(\frac{1}{3}\right) \approx 70\dec 53^\circ.
\]

\subsection*{Défi bonus}

a) En écrivant $\dfrac{\sin(\alpha)}{\cos(\alpha)} = \dfrac{\text{opp}/\text{hyp}}{\text{adj}/\text{hyp}} = \dfrac{\text{opp}}{\text{hyp}} \times \dfrac{\text{hyp}}{\text{adj}} = \dfrac{\text{opp}}{\text{adj}} = \tan(\alpha)$.\\
b) $0\mbox{,}6^2 + \cos^2(\alpha) = 1 \implies 0\mbox{,}36 + \cos^2(\alpha) = 1 \implies \cos^2(\alpha) = 0\mbox{,}64 \implies \cos(\alpha) = \sqrt{0\mbox{,}64} = 0\mbox{,}8$.

\end{document}